\scanpage{7}
\mspage{4}

2. Préliminaires heuristiques sur les dérivateurs.

On se donne une sous-catégorie plei\deleted{\unsure{n}}\unsure{ne} $\mathrm{Dia}$ (comme ``diagrammes'') de $\mathrm{Cat}$. Ça peut être tout $\mathrm{Cat}$, ou plus petit. Mais on supposera le plus souvent que $\mathrm{Dia}$ contient au moins les ``catégories ordonnées'' (correspondant \unsure{aux} \unsure{ens.|ens} ordonnés, i.e. telles que a) tout morphisme est \unsure{mono} et b) tout iso est une identité). Par abus de langage, on identifie un ens. ordonné avec la cat. correspondante, (\unsure{et} une cat. au topos qu'elle définit ?).

% slanted
\margin{\uline{NB} Il faut des conditions de stabilité pour $\mathrm{Dia}$, au moins la stabilité pour $I\times J$ et $I\unsure{\amalg|\sqcup}J$ -- en fait, nettement plus\ldots}

On considère $\mathrm{Dia}$ comme une 2-\uline{catégorie}.

Un \uline{dérivateur} \inserted{$\unsure{\mathbb D|D}$} (sur $\mathrm{Dia}$, comme ``\uline{domaine}'' du dérivateur) est un 2-foncteur contravariant de $\mathrm{Dia}$ dans $\mathrm{Cat}$, satisfaisant certaines conditions. Donc \unsure{A|à}

\[
\mathbb D:\mathrm{Dia}^{\unsure{\circ|o}}\longrightarrow\mathrm{Cat}
\]

Si $I\in\mathrm{Ob}_{\unsure{0|\circ}}\,\mathrm{Dia}$ est une catégorie (ou un ``type de diagramme'') du domaine de $D$, on doit donc avoir une catégorie $\mathbb D(I)$. Les objets de $D(I)$ s'appelleront les objets-coefficients (ou \uline{\unsure{coéfficients|coefficients}}) de type $D$, au dessus de $I$. Pour

\[
f:I\unsure{-|\to}J
\]

une 1-flèche de $\mathrm{Dia}$, on doit avoir un foncteur

\[
f^{*}_{\mathbb D}\ \text{ou}\ f^{*}:\mathbb D(J)\longrightarrow\mathbb D(I),
\]

dit ``\uline{foncteur image inverse}''. Si $F$ est un coeff. de type $\mathbb D$ sur $J$, son ``\uline{image inverse}'' est donc un coeff. de type $\mathbb D$ sur $I$. On suppose \inserted{(pour simplifier)} que le 2-foncteur $\mathbb D$ est \uline{strict}, donc qu'on a transitivité stricte

\[
(gf)^{*}=f^{*}g^{\unsure{*|\prime}}\qquad\text{si}\quad \underset{\text{dans }\mathrm{Dia}}{I\overset{f}{-}J\overset{g}{\longrightarrow}K}
\]

% archive sheet number: 6
